\documentclass[]{../../../../tex/classes/styledArticle}
\usepackage{../../../../tex/packages/hebrewSupport}
\usepackage{../../../../tex/packages/mathShortcuts}
\usepackage{../../../../tex/packages/theoremsSupport}

\author{שחר פרץ}
\title{תרגול לצוציקים של שנה א'}
\begin{document}
	\maketitle
	\subsection*{דוגמה לחזרה על לכסון}
	\exe{הוכיחו ש־$\sof A > \az$
		\[ A := \ccb{f \co \N \to \N \co \forall n \in \Neven \co f(n) = 0} \]}
	\begin{proof}[\textbf{לא מלאה}, חסרים פרטים]
		\begin{itemize}
			\item $\bm{\sof A  \ge \az}$\textbf{: }נתבונן בפונקציה החח''ע הבאה: 
			\[ f \co \N \to A \quad \quad f = \lambda n \in \N. \lambda m \in \N\co \begin{cases}
				n & m \in \Nodd \\
				0 & n \in \Neven
			\end{cases} \]
			יש להוכיח שהיא (א) חח''ע ו(ב) מוגדרת היטב. הראנו זאת בתרגול. 
			\item $\bm{\sof A \neq \con}$\textbf{: }נניח בשלילה שקיימת $F \co \N \to A$ חח''ע ועל. נתבונן בפונקציה הבאה: 
			\[ f \co \N \to \N \quad\quad f = \lambda n \in \N\co \begin{cases}
				F\cl{\frac{n - 1}{2}}(n) + 1 & n \in \Nodd \\
				0 & n \in \Neven
			\end{cases} \]
			יש להראות ש$f \notin \Img F$ וגם $ f \in A$, עשינו זאת בתרגול. זו סתירה לכך ש־$f$ על. 
		\end{itemize}\envendproof
	\end{proof}
	
	תרגילים דומים: \exe{הוכיחו ש־: 
			\[ A = \ps(\N \times \N) \setminus (\N \to \N) \]
			מקיימת $\sof{A} > \az$. }\exe{$C = \{f \co \N \to \N\mid f[\Neven] \cap f[\Nodd] = \varnothing \}$. הוכיחו ש־$\sof C > \az$ באמצעות לכסון. }
	
	התרגילים לעיל דומים מאוד, אבל אני ממליץ לכם לנסות לפתור אותם במידה ואתם לא חושבים שתוכלו לכתוב את ההוכחה מסודר לבד. 
	
	\subsection*{שאלה ראשונה}
	\exe{הוכיחו ש־: 
		\[ A = \ps(\N \times \N) \setminus (\N \to \N) \]
		מקיימת $\sof{A} > \az$. }
	\begin{itemize}
		\item גדול חלש: 
		\[ \lambda n \in \N\co \ccb{\la n, 0 \ra} \]
		יש להראות שהפונקציה לעיל חח''ע וטווחה מוגדר היטב. 
		\item לא שווה: בשלילה $F(n) \co \N \to A$ חח''ע ועל, נגדיר $F(n) = R_n$ ונבחר: 
		\[ R = \{\la 0, 1 \ra , \la 0, 0 \ra\} \cup \bigcup_{n \in \N_+}\{\la a, b \ra \in \N^{2}\co a = n \land \lnot aR_{n - 1}b\} \]
		ראשית, נבחין ש־$R \in A$ כי $R \subset \N \times \N$ ולכן $R \in \ps(\N \times \N)$, אך $R \notin \N\to \N$ מכיוון שהיא לא חד־ערכית (מתקיים $0R1 \land 0R 0$). 
		
		מהיות $f$ על, בהכרח $\exists n\in \N \co F(n) = R_n = R$. אזי בפרט $(n + 1)R_nb \iff (n + 1)Rb$ ובפרט עבור $b = 1$. נפרק למקרים. 
		\begin{itemize}
			\item אם $\la n + 1, 1 \ra \in R_n$, אז $(n + 1)R_{n + 1 - 1}b$, כלומר$\la n + 1, b \ra \notin R$ מהגדרה (וכי $n + 1 \in \N_+$) וזו סתירה לשוויון קבוצות. 
			\item אם $\la n + 1, 1 \ra \notin R_n$, אז באופן דומה $\la n + 1, 1 \ra \in R$ וזו סתירה. 
		\end{itemize}
		סה''כ $R_n \neq R$ וזו סתירה לכך ש־$F(n) = R$, כנדרש. 
	\end{itemize}
	
	\subsection*{שאלה שנייה}
	23CA במבוא לתקב''צ. 
	
	\exe{תהי $B$ אינסופית, ותהי $A$ כך ש־$\sof a \le \sof B$. הוכיחו שקיימת חלוקה $(C_i)_{i \in I}$  כך ש־$\sof{C_i} = \sof B$ ו־$\sof I = \sof A$. }
	\textit{הערה. }הכוונה ב־$(C_i)_{i \in I}$, היא למצוא $I$ מעוצמה $\sof A$ ולמצוא פונקציה $F \co I \to \ps(B)$, ונסמן $C_i := F(i)$. 
	
	\begin{proof}[הוכחה \textbf{לא מלאה, חסרים פרטים}]
		ידוע $\sof A \cdot \sof B = \sof B$ מחשבון עוצמות, כי ממונוטוניות $\sof A \cdot \sof B \ge \sof B$ אך גם $\sof B \cdot \sof A \le \sof B \cdot \sof B = \sof B$. לכן קיים זיווג $f \co I \times B \to B$. נגדיר:
		\[ C_i = f[\{i\} \times B] \]
		עבור $i \in I = A$. 
		עתה יש להוכיח שזו חלוקה המקיימת את הנדרש. 
		\begin{itemize}
			\item ידוע $I = A$ כלומר $\sof I = \sof A$
			\item ידוע $f$ זיווג לכן $B \ni b \mapsto f(\la i, b\ra) \in C_i$ זיווג. מכאן ש־$\sof B = \sof {C_i}$
			\item זו אכן חלוקה: 
			\begin{itemize}
				\item מתקיים ש־$C_i \neq \varnothing$ כי $\sof{C_i} = \sof B \neq 0$. 
				\item מתקיים ש־$C_i = C_j \iff i = j$ (זרים בזוגות) כי $f$ חח''ע. 
				\item מתקיים ש־$\bigcup C_i = B$ כי: 
				\[ \bigcup_{i \in I} C_i = \bigcup_{i \in I}f[\{i\} \times B] = f \csb{\bigcup_{i \in I} \{i\}\times B} = f\csb{A \times B} = B \]
			\end{itemize}
		\end{itemize}
	\end{proof}
	
	\subsection*{הערות נוספות}
	\begin{itemize}
		\item מי שרוצה ללמוד LyX או \en{\LaTeX}\ מוזמן לפנות אלי במייל או בטלפון. אגב, נטלי מעבירה בפקולטה הרצאות על ליך מדי פעם. אתם יכולים לשאול אותה מתי זה קורה ולהגיע לאחת מהן. 
		\item שאלו אותי האם שכאשר רוצים להוכיח ש־$\sof A > \sof B$, האם אפשר להוכיח שלא קיימת פונקציה $f \co B \to A$ שהיא על, במקום להוכיח שקיימת $f \co B \to A$ חח''ע ושלא קיימת $F \co B \to A$ שהיא זיווג. התשובה היא כן: זה הקונטראפוזיטיב של משפט שהוכחתם בכיתה, בהינתן אקסיומת הבחירה. 
		\begin{proof}
			נניח שלא קיימת $f \co B \to A$ על. נניח בשלילה $\sof A \ge \sof B$, ואז מאקסיומת הבחירה הוכחתם בכיתה שקיימת $f \co B \to A$ על, וזו סתירה. 
		\end{proof}
		כאשר מבקשים מכם ''להוכיח באמצעות לכסון`` הכוונה היא להראות קיום פונקציה חח''ע ושלא קיים זיווג, ולא להראות שלא קיימת פונקציה על כמו שמתואר לעיל. 
		\item המבחנים של בדידה 1 לפני 2024 (לא כולל) שנמצאים בדרייב די קלים, ולא הייתי מצפה שהמבחן שלכם יראה ככה. הסיבה לכך היא שלפני 2024, הקורס היה שילוב של בדידה 1 עם בדידה 2, ובאופן טבעי השאלות שנגעו בחומר של בדידה 1 היו יותר קלות. הייתי שוקל להתחיל ללמוד מהמבחנים של בדידה 1, ואחת שאתם מרגישים בנוח עם החומר לתרגל מהמבחנים של ''מבוא לתקב''צ`` (קורס של בית הספר למתמטיקה). שימו לב שבמבוא לתקב''צ יש חומר ש\textbf{לא} נמצא בבדידה 1, והוא: 
		\begin{itemize}
			\item שרשראות, הלמה של צורן
			\item משפט הסדר הטוב
			\item איזומורפיזם של יחסי סדר
		\end{itemize}
		באופן אישי הייתי ממליץ לכם ללמוד על הלמה של צורן ועל שרשראות לידע כללי (זה יועיל לכם בקורס מאוחרים יותר כמו טופולוגיה) אבל זה לא קריטי למבחן שלכם. 
	\end{itemize}
	
	
	
	\ndoc
\end{document}